\documentclass[a4paper,11pt]{article}

\usepackage[french]{babel}
\frenchbsetup{StandardLists=true}
\usepackage[utf8]{inputenc}
\usepackage[T1]{fontenc}
\usepackage[top=2.5cm, bottom=2.5cm, left=2.5cm, right=2.5cm]{geometry}
\usepackage{lmodern}
\usepackage{xcolor}
\usepackage{hyperref}
\usepackage{float}
\usepackage{graphicx}
\usepackage{array}
\usepackage{tabularx}
\usepackage{longtable}
\usepackage{booktabs}
\usepackage{enumitem}
\usepackage{fancyhdr}
\usepackage{titlesec}
\usepackage{caption}
\usepackage{listings}
\usepackage[most]{tcolorbox}
\usepackage{mdframed}
\usepackage{makecell}

%================ COLORS =================
\definecolor{bleu}{RGB}{0,76,153}
\definecolor{bleuclair}{RGB}{0,102,204}
\definecolor{grislight}{RGB}{245,245,244}
\definecolor{grisborder}{RGB}{210,210,210}
\definecolor{orange}{RGB}{200,80,0}
\definecolor{vert}{RGB}{0,120,60}
\definecolor{rouge}{RGB}{180,20,20}
\definecolor{jaune}{RGB}{180,130,0}

\hypersetup{
    colorlinks=true,
    linkcolor=bleu,
    urlcolor=bleuclair,
    citecolor=bleu
}

%================ HEADER / FOOTER =================
\pagestyle{fancy}
\fancyhf{}
\fancyhead[L]{\small\color{bleu}\textbf{Infrastructure HCI Open Source, Documentation Generale}}
\fancyhead[R]{\small\color{bleu}Juin 2026}
\fancyfoot[C]{\small\thepage}
\renewcommand{\headrulewidth}{0.4pt}
\renewcommand{\footrulewidth}{0pt}

%================ SECTION TITLES =================
\titleformat{\section}
  {\normalfont\Large\bfseries\color{bleu}}
  {\thesection}{1em}{}[\vspace{-0.2em}\rule{\linewidth}{0.4pt}]

\titleformat{\subsection}
  {\normalfont\large\bfseries\color{bleuclair}}
  {\thesubsection}{1em}{}

\titleformat{\subsubsection}
  {\normalfont\normalsize\bfseries\color{bleu!70!black}}
  {\thesubsubsection}{1em}{}

%================ CODE BOX =================
\lstdefinestyle{shellstyle}{
    backgroundcolor=\color{grislight},
    frame=single,
    rulecolor=\color{grisborder},
    basicstyle=\ttfamily\small,
    keywordstyle=\color{bleu}\bfseries,
    commentstyle=\color{vert},
    stringstyle=\color{rouge},
    breaklines=true,
    showstringspaces=false,
    tabsize=2
}

%================ WARNING / INFO BOXES =================
\newenvironment{infobox}{%
  \begin{mdframed}[linecolor=bleu,backgroundcolor=bleu!5,linewidth=1pt,
    leftmargin=0pt,rightmargin=0pt,innertopmargin=6pt,innerbottommargin=6pt]
}{%
  \end{mdframed}
}

\newenvironment{warnbox}{%
  \begin{mdframed}[linecolor=orange,backgroundcolor=orange!8,linewidth=1pt,
    leftmargin=0pt,rightmargin=0pt,innertopmargin=6pt,innerbottommargin=6pt]
}{%
  \end{mdframed}
}

\newenvironment{successbox}{%
  \begin{mdframed}[linecolor=vert,backgroundcolor=vert!7,linewidth=1pt,
    leftmargin=0pt,rightmargin=0pt,innertopmargin=6pt,innerbottommargin=6pt]
}{%
  \end{mdframed}
}

%================ TABLE STYLE =================
\newcolumntype{L}[1]{>{\raggedright\arraybackslash}p{#1}}
\newcolumntype{C}[1]{>{\centering\arraybackslash}p{#1}}

%================ DOCUMENT =================
\begin{document}

%----------- PAGE DE TITRE -----------
\begin{titlepage}
  \centering
  \vspace*{2cm}
  {\color{bleu}\rule{\linewidth}{1.5pt}}\vspace{0.4cm}

  {\Huge\bfseries\color{bleu} Infrastructure Hyperconvergee\\[0.3em]
  Open Source}\par\vspace{0.5cm}
  {\color{bleu}\rule{\linewidth}{0.6pt}}\par\vspace{0.8cm}

  {\LARGE Documentation Generale d'Exploitation}\par\vspace{0.4cm}
  {\large\color{gray} Etat de l'infrastructure, Juin 2026}\par\vspace{2cm}

  \begin{tabular}{ll}
    \textbf{Realise par :} & Bouzara Zakaria \& Ilyes \\[0.4em]
    \textbf{Contexte :}    & Projet de Fin d'Etudes, Master 2 RSD \\[0.4em]
    \textbf{Encadrement :} & USTHB \\[0.4em]
    \textbf{Date :}        & Juin 2026 \\
  \end{tabular}

  \vfill
  {\color{bleu}\rule{\linewidth}{1.5pt}}
\end{titlepage}

\tableofcontents
\newpage

%============================================================
\section{Presentation generale du projet}
%============================================================

Ce document constitue la documentation technique d'exploitation de l'infrastructure
hyperconvergee deployee dans le cadre du projet de fin d'etudes. Il decrit l'etat
reel de chaque composant~: ce qui fonctionne, ce qui presente des limitations,
et les points a completer ou corriger.

L'objectif du projet est de concevoir et deployer une infrastructure \textit{cloud prive}
entierement basee sur des logiciels libres, reproduisant les fonctionnalites essentielles
d'une plateforme d'entreprise~: hyperconvergence, haute disponibilite, ordonnancement
des ressources, segmentation reseau, supervision et automatisation.

%============================================================
\section{Infrastructure physique}
%============================================================

\subsection{Materiel disponible}

L'infrastructure repose sur trois serveurs physiques \textbf{Dell PowerEdge R720},
connectes entre eux via un commutateur physique et relies a Internet a travers
un ordinateur portable faisant office de passerelle NAT.

\begin{table}[H]
  \caption{Serveurs physiques du cluster}
  \centering
  \begin{tabular}{L{2.5cm}L{3.5cm}L{3.5cm}L{4cm}}
    \toprule
    \textbf{Hote} & \textbf{IP de gestion} & \textbf{Role} & \textbf{Remarques} \\
    \midrule
    pve-01 & 192.168.10.100 & Noeud Proxmox & Noeud primaire \\
    pve-02 & 192.168.10.213 & Noeud Proxmox & MGR Ceph actif \\
    pve-03 & 192.168.10.84  & Noeud Proxmox & MGR Ceph standby \\
    Laptop  & 192.168.10.36  & Passerelle NAT / admin & Acces internet \\
    \bottomrule
  \end{tabular}
\end{table}

\subsection{Acces SSH et identifiants}

Tous les composants de l'infrastructure sont accessibles par SSH. Le tableau ci-dessous
recapitule les identifiants pour chaque composant.

\begin{table}[H]
  \caption{Identifiants d'acces a l'infrastructure}
  \centering
  \begin{tabular}{L{4.5cm}L{3.2cm}L{4.5cm}}
    \toprule
    \textbf{Composant} & \textbf{Identifiants} & \textbf{Remarques} \\
    \midrule
    Serveurs Proxmox (pve-01/02/03) & \texttt{root:ZakiIlyes} & SSH + interface web \\
    VMs et conteneurs (Sona-Web, CLBS, monitoring\ldots) & \texttt{root:ZakiIlyes} & SSH direct \\
    GitLab & \texttt{root:ZakiIlyes7} & Interface web + API \\
    Semaphore UI & \texttt{admin:admin} & Interface playbooks Ansible \\
    \bottomrule
  \end{tabular}
\end{table}

\begin{infobox}
  \textbf{Connexion SSH aux noeuds Proxmox~:}\\
  \texttt{ssh root@192.168.10.100} \quad (pve-01)\\
  \texttt{ssh root@192.168.10.213} \quad (pve-02)\\
  \texttt{ssh root@192.168.10.84}  \quad (pve-03)
\end{infobox}

%============================================================
\section{Cluster Proxmox VE}
%============================================================

\subsection{Etat general}

Les trois noeuds forment un cluster Proxmox VE pleinement operationnel.
Corosync assure la gestion du quorum et du \textit{heartbeat} inter-noeuds.
Le cluster est accessible depuis l'interface web sur le port \texttt{8006}
de n'importe quel noeud.

\begin{successbox}
  \textbf{Statut~: Fonctionnel} \\
  Cluster a 3 noeuds, quorum maintenu, API Proxmox disponible, migrations a chaud operationnelles.
\end{successbox}

Les commandes de verification de l'etat du cluster~:

\begin{lstlisting}[style=shellstyle]
pvecm status        # etat du cluster et du quorum
pvecm nodes         # liste des noeuds
ceph -s             # etat du stockage distribue
\end{lstlisting}

%============================================================
\section{Reseau: OVS et VxLAN}
%============================================================

\subsection{Architecture reseau}

Le reseau de l'infrastructure est entierement gere par \textbf{Open vSwitch (OVS)}.
Neuf reseaux logiques distincts ont ete crees sous forme de \textbf{VxLAN}
(et non de VLAN traditionnels).

\begin{warnbox}
  \textbf{Important: Limitation de gestion reseau~:}\\
  L'infrastructure utilise des tunnels \textbf{VxLAN} (protocole RFC~7348, encapsulation UDP
  sur le port~4789) et non des VLAN 802.1Q classiques. L'interface web de Proxmox VE
  \textbf{ne supporte pas nativement la configuration VxLAN}~: toute modification
  des parametres reseau (ajout d'interface, modification d'adresses, changement de
  configuration OVS) doit imperativement etre effectuee \textbf{en ligne de commande}
  via \texttt{/etc/network/interfaces}, et non depuis l'interface graphique.
  Toute tentative de modification via l'interface web risque de supprimer la
  configuration VxLAN existante.
\end{warnbox}

\subsection{Plan d'adressage VxLAN}

\begin{table}[H]
  \caption{Reseaux VxLAN de l'infrastructure}
  \centering
  \begin{tabular}{C{1.5cm}L{3.5cm}L{3.5cm}L{4cm}}
    \toprule
    \textbf{VxLAN} & \textbf{Nom} & \textbf{Sous-reseau} & \textbf{Usage} \\
    \midrule
    \makecell{10\\~} &
    \makecell[l]{Management\\~} &
    \makecell[l]{192.168.10.0/24\\~} &
    \makecell[l]{Web UI, SSH, API Proxmox\\~} \\
    \makecell{20\\~} &
    \makecell[l]{Corosync\\~} &
    \makecell[l]{192.168.20.0/24\\~} &
    \makecell[l]{Heartbeat, quorum\\~} \\
    \makecell{30\\~} &
    \makecell[l]{Ceph Cluster\\~} &
    \makecell[l]{10.10.30.0/24\\~} &
    \makecell[l]{Replication OSD interne\\~} \\
    \makecell{40\\~} &
    \makecell[l]{Ceph Public\\~} &
    \makecell[l]{10.10.40.0/24\\~} &
    \makecell[l]{Acces RBD depuis les VMs\\~} \\
    \makecell{50\\~} &
    \makecell[l]{Live Migration\\~} &
    \makecell[l]{10.10.50.0/24\\~} &
    \makecell[l]{Migrations a chaud (CLBS)\\~} \\
    \makecell{60\\~} &
    \makecell[l]{Production\\~} &
    \makecell[l]{172.16.60.0/24\\~} &
    \makecell[l]{VMs applicatives Sona-Web\\~} \\
    \makecell{70\\~} &
    \makecell[l]{Dev / Test\\~} &
    \makecell[l]{172.16.70.0/24\\~} &
    \makecell[l]{Environnement de test\\~} \\
    \makecell{80\\~} &
    \makecell[l]{DMZ\\~} &
    \makecell[l]{172.16.80.0/24\\~} &
    \makecell[l]{HAProxy, services exposes\\~} \\
    \makecell{90\\~} &
    \makecell[l]{Backup\\~} &
    \makecell[l]{10.10.90.0/24\\~} &
    \makecell[l]{Proxmox Backup Server\\~} \\
    \bottomrule
  \end{tabular}
\end{table}

\subsection{Etat du reseau}

\begin{successbox}
  \textbf{Statut~: Fonctionnel} \\
  OVS actif sur les trois noeuds, tous les tunnels VxLAN etablis,
  connectivite inter-noeuds verifiee.
\end{successbox}

%============================================================
\section{Stockage distribue: Ceph}
%============================================================

\subsection{Configuration Ceph}

Le cluster Ceph est deploye sur les trois noeuds physiques selon la configuration suivante~:

\begin{itemize}[noitemsep]
  \item \textbf{3 Monitors (MON)}~: un par noeud, assurant le quorum Ceph
  \item \textbf{2 Managers (MGR)}~: pve-02 (actif) et pve-03 (standby)
  \item \textbf{3 OSD}~: un disque dedie par noeud (facteur de replication~: 3)
  \item \textbf{Pool}~: \texttt{vmpool}, replication~x3, \texttt{pg\_autoscale} active
  \item Reseau cluster~: VxLAN~30 (\texttt{10.10.30.0/24})
  \item Reseau public~: VxLAN~40 (\texttt{10.10.40.0/24})
\end{itemize}

\begin{successbox}
  \textbf{Statut~: Fonctionnel} \\
  Ceph en etat \texttt{HEALTH\_OK}, tous les OSD \texttt{up+in}, replication active.
  Le pool \texttt{vmpool} est le stockage par defaut pour les disques de VMs.
\end{successbox}

Commandes de verification~:

\begin{lstlisting}[style=shellstyle]
ceph -s              # etat global (HEALTH_OK attendu)
ceph osd tree        # arborescence des OSD
ceph df              # utilisation du stockage
\end{lstlisting}

%============================================================
\section{Haute disponibilite (HA)}
%============================================================

\subsection{Configuration}

La haute disponibilite est geree par le gestionnaire HA integre de Proxmox VE.
Les VMs critiques (\textit{Sona-Web}, \textit{CLBS}) sont enregistrees dans le gestionnaire HA
avec une politique de redemarrage automatique sur un noeud sain en cas de defaillance.

\subsection{Test de basculement}

Le basculement HA a ete teste en simulant la perte d'un noeud via le blocage du
port Corosync (UDP~5405). Le test a confirme que~:

\begin{itemize}[noitemsep]
  \item Le CRM detecte la perte du noeud apres environ 30 secondes
  \item Les VMs sont redemarrees sur les noeuds survivants
  \item Le watchdog (\textit{softdog}) assure le \textit{fencing} du noeud defaillant
\end{itemize}

\begin{warnbox}
  \textbf{Prerequis HA~:} Sans watchdog/fencing configure, Proxmox hesite a effectuer
  le basculement pour eviter le \textit{split-brain}. La configuration du watchdog
  \texttt{softdog} est indispensable au bon fonctionnement du HA.
\end{warnbox}

\begin{successbox}
  \textbf{Statut~: Fonctionnel} \\
  HA valide par test de basculement. Redemarrage automatique des VMs confirme.
\end{successbox}

%============================================================
\section{Routage et pare-feu: OPNsense}
%============================================================

\subsection{Role d'OPNsense}

Une VM OPNsense fait office de routeur et de pare-feu inter-zones pour les VMs
applicatives. Les noeuds Proxmox eux-memes conservent le laptop (192.168.10.36)
comme passerelle afin de decoupler la disponibilite du cluster de l'etat de la
VM OPNsense.

\subsection{Etat des fonctionnalites OPNsense}

\subsubsection{NAT et routage}

\begin{successbox}
  \textbf{NAT~: Fonctionnel}: Le NAT sortant est actif (mode automatique).
  Les VMs des zones PROD, TEST et DMZ acce\-dent a Internet via la VM OPNsense.\\[0.3em]
  \textbf{Routage inter-zones~: Fonctionnel}: Le routage entre les zones PROD (172.16.60.0/24),
  TEST (172.16.70.0/24) et DMZ (172.16.80.0/24) est operationnel.
\end{successbox}

\subsubsection{DHCP}

\begin{warnbox}
  \textbf{DHCP~: Partiellement fonctionnel: Conflit intermittent.}\\
  Le serveur DHCP d'OPNsense entre en conflit avec le commutateur physique
  qui dispose lui-meme d'un serveur DHCP actif. Il s'agit d'une \textit{course DHCP}
  (\textit{DHCP race}) entre les deux serveurs~: selon lequel repond en premier,
  les VMs peuvent recevoir une adresse incorrecte ou ne pas en recevoir du tout.
  Les adresses IP statiques sont recommandees pour contourner ce probleme.
\end{warnbox}

\subsubsection{Pare-feu inter-zones (PROD / TEST / DMZ)}

\begin{warnbox}
  \textbf{Statut~: Partiellement valide.}\\
  Les regles de segmentation entre les zones PROD, TEST et DMZ sont configurees
  et fonctionnelles dans les cas standards. Les cas limites (ports applicatifs
  specifiques, protocoles non standard) n'ont pas ete exhaustivement testes.
\end{warnbox}

\subsubsection{DNS}

\begin{warnbox}
  \textbf{DNS~: Active mais non configure.}\\
  Le serveur DNS d'OPNsense est active mais aucune regle ni zone DNS personnalisee
  n'a ete definie. Il transfere les requetes vers les serveurs upstream par defaut
  (8.8.8.8).
\end{warnbox}

\subsubsection{Desactivation temporaire du pare-feu OPNsense}

Lors des phases de configuration ou de debogage, il est parfois necessaire de
desactiver temporairement le pare-feu d'OPNsense pour acceder a son interface
web depuis l'adresse IP de la VM. La commande a executer depuis la console
OPNsense (ou via SSH) est~:

\begin{lstlisting}[style=shellstyle]
pfctl -d    # desactive le packet filter (pare-feu) d'OPNsense
\end{lstlisting}

\textbf{Attention~:} Cette commande desactive entierement le filtrage de paquets.
Elle doit rester temporaire et etre annulee avec \texttt{pfctl -e} apres configuration.

%============================================================
\section{Pare-feu Proxmox VE}
%============================================================

\subsection{Architecture et etat}

Le pare-feu integre de Proxmox VE a ete configure et active sur les trois noeuds.
Il s'appuie sur \texttt{iptables} et s'organise en trois niveaux~: Datacenter, Noeud, et VM.

\begin{successbox}
  \textbf{Statut~: Fonctionnel et active.}\\
  Le pare-feu est en mode \texttt{DROP} en entree. Les groupes de securite
  (Corosync, VxLAN, Ceph, migration a chaud, administration) sont definis
  au niveau Datacenter et propage aux trois noeuds.
\end{successbox}

\subsection{Limitation~: ports applicatifs non declares}

\begin{warnbox}
  \textbf{Ports de developpement non ouverts.}\\
  Les ports utilises par les services en developpement (par exemple~: \texttt{8000},
  \texttt{5432} (PostgreSQL depuis l'exterieur), \texttt{5423} et autres ports applicatifs
  non standards) ne sont pas definis dans les regles du pare-feu Proxmox. Ces ports doivent
  etre ouverts \textbf{a la demande}, au moment ou le service correspondant en a besoin.\\[0.3em]
  Pour ouvrir un port sur une VM~: \textit{Proxmox Web UI} puis
  selectionner la VM, puis \textit{Firewall}, puis \textit{Rules}, puis bouton \textit{Add}.
\end{warnbox}

\subsection{Ports cles configures}

\begin{table}[H]
  \caption{Principaux ports ouverts dans le pare-feu Proxmox}
  \centering
  \begin{tabular}{C{2.5cm}C{2cm}L{4cm}L{4cm}}
    \toprule
    \textbf{Port(s)} & \textbf{Proto} & \textbf{Service} & \textbf{Reseau source} \\
    \midrule
    22            & TCP & SSH administration       & VxLAN~10 (192.168.10.0/24) \\
    8006          & TCP & Web UI + API Proxmox     & VxLAN~10 + noeuds pairs \\
    4789          & UDP & Tunnels VxLAN (CRITIQUE) & Noeuds physiques \\
    5404--5405    & UDP & Corosync heartbeat       & VxLAN~20 \\
    3300, 6789    & TCP & Ceph MON v2/v1           & VxLAN~30, 40 \\
    6800--7300    & TCP & Ceph OSD                 & VxLAN~30, 40 \\
    60000--60050  & TCP & QEMU live migration      & VxLAN~50 \\
    80, 443       & TCP & HTTP/HTTPS (HAProxy)     & Toutes origines \\
    3000          & TCP & Grafana                  & VxLAN~10 \\
    8086          & TCP & InfluxDB API             & VxLAN~10 \\
    \bottomrule
  \end{tabular}
\end{table}

%============================================================
\section{Ordonnancement dynamique: CLBS}
%============================================================

\subsection{Description}

Le \textit{Custom Load Balancing Service} (CLBS) est un service ecrit en \textbf{Go}
qui surveille en temps reel l'utilisation CPU et memoire des noeuds Proxmox
via l'API REST, et declenche des migrations a chaud (\textit{live migration})
pour reequilibrer la charge entre les noeuds.

\begin{itemize}[noitemsep]
  \item Seuil de declenchement~: 80\,\% CPU ou memoire sur un noeud
  \item Seuil bas~: 30\,\% (noeud cible des migrations)
  \item Periode de refroidissement~: 60 secondes entre deux migrations
  \item Interface web de supervision~: \textit{Sona-Web} affiche les migrations
        effectuees et l'etat du cluster en temps reel
\end{itemize}

\begin{successbox}
  \textbf{Statut~: Fonctionnel.}\\
  Le service CLBS est deploye, les migrations automatiques ont ete validees.
  L'interface web Sona-Web reflete l'etat des migrations en temps reel.
\end{successbox}

Les sources se trouvent dans \texttt{codes/CLBS-golang/}~: \texttt{main.go},
\texttt{metrics.go}, \texttt{store.go}.

\subsection{Perspective: CRS natif de Proxmox}

Les versions recentes de Proxmox VE integrent un mecanisme natif d'equilibrage
de charge appele \textbf{CRS} (\textit{Cluster Resource Scheduler}). Il est preferable
au CLBS sur plusieurs points~:

\begin{itemize}[noitemsep]
  \item Le CRS fonctionne comme un service systeme sur les noeuds Proxmox eux-memes,
        et non dans une VM qui peut etre indisponible au moment d'un basculement HA.
  \item Il utilise le quorum Corosync pour la communication et la synchronisation
        entre noeuds, ce qui le rend plus reactif et plus robuste que le mecanisme
        de secours par HA d'une VM.
  \item La reintegration dans l'ecosysteme Proxmox est native, sans API externe.
\end{itemize}

\begin{infobox}
  \textbf{Recommandation~:} Il est envisageable de creer un \textit{playbook} Ansible
  pour mettre a jour la version de Proxmox sur les trois noeuds, puis de basculer
  du CLBS vers le CRS. Cette evolution permettrait de supprimer la VM CLBS et de
  beneficier d'un ordonnancement plus fiable, directement integre dans le cluster.
\end{infobox}

%============================================================
\section{Application Sona-Web}
%============================================================

\subsection{Architecture applicative}

Sona-Web est l'application de demonstration deployee sur l'infrastructure.
Elle est constituee de quatre VMs applicatives~:

\begin{table}[H]
  \caption{VMs de l'application Sona-Web}
  \centering
  \begin{tabular}{L{2.5cm}L{3cm}C{2.5cm}L{4cm}}
    \toprule
    \textbf{VM} & \textbf{Role} & \textbf{Reseau} & \textbf{Remarques} \\
    \midrule
    haproxy     & Repartiteur de charge  & VxLAN 60 + 80 & Equilibrage round-robin \\
    flask-app-1 & Serveur Flask (Python) & VxLAN 60      & Application backend \\
    flask-app-2 & Serveur Flask (Python) & VxLAN 60      & Replique pour HA \\
    postgresql  & Base de donnees        & VxLAN 60      & Disque sur Ceph RBD \\
    \bottomrule
  \end{tabular}
\end{table}

\begin{successbox}
  \textbf{Statut~: Fonctionnel.}\\
  L'application est accessible via HAProxy. La repartition de charge est active.
  La base de donnees PostgreSQL est accessible depuis les deux instances Flask.
\end{successbox}

Les sources se trouvent dans \texttt{codes/sona-web/}~: \texttt{app.py} (Flask),
\texttt{haproxy.cfg}, schema de la base de donnees.

%============================================================
\section{Supervision: Conteneur de monitoring}
%============================================================

\subsection{Stack de supervision}

La supervision de l'infrastructure est assuree par trois services regroupes dans un
\textbf{conteneur LXC dedie} (CT~102) deploye sur Proxmox~:

\begin{itemize}[noitemsep]
  \item \textbf{InfluxDB}~: collecte des metriques Proxmox (CPU, memoire, reseau par noeud et par VM),
        envoyes par les noeuds via \textit{Metric Server} (\texttt{Datacenter > Metric Server > InfluxDB})
  \item \textbf{Prometheus}~: active via le module \textit{Prometheus MGR} de Ceph pour l'exposition
        des metriques Ceph
  \item \textbf{Loki}~: service de centralisation des logs (deploye mais non alimente en sources de logs)
  \item \textbf{Grafana}~: visualisation des metriques, accessible sur le port~\texttt{3000}
\end{itemize}

\begin{successbox}
  \textbf{Statut~: Fonctionnel.}\\
  Grafana operationnel, metriques des noeuds et VMs visibles.
  Prometheus MGR de Ceph actif.
\end{successbox}

\subsection{Acces au conteneur de monitoring}

\begin{warnbox}
  \textbf{Acces depuis l'interface Proxmox non disponible.}\\
  Le conteneur de monitoring (CT~102) ne dispose pas de l'agent QEMU guest.
  La console integree de l'interface web Proxmox ne fonctionne pas pour ce conteneur.
  Deux methodes d'acces sont disponibles~:
\end{warnbox}

\begin{enumerate}[noitemsep]
  \item \textbf{SSH direct} (recommande)~:
\begin{lstlisting}[style=shellstyle]
ssh root@192.168.10.99   # mot de passe : ZakiIlyes
\end{lstlisting}
  \item \textbf{Entree directe depuis le noeud hote} (si SSH indisponible)~:
\begin{lstlisting}[style=shellstyle]
pct enter 102   # depuis pve-02 (ou le noeud qui heberge CT 102 apres HA)
\end{lstlisting}
\end{enumerate}

Le conteneur est normalement heberge sur \textbf{pve-02}. S'il a ete deplace par le
mecanisme HA, verifier sur quel noeud il se trouve depuis l'interface web Proxmox
avant d'utiliser \texttt{pct enter}.

\subsection{Fichiers de configuration}

\begin{table}[H]
  \caption{Fichiers de configuration des services de supervision (CT~102)}
  \centering
  \begin{tabular}{L{3cm}L{5.5cm}L{5cm}}
    \toprule
    \textbf{Service} & \textbf{Fichier de configuration} & \textbf{Remarques} \\
    \midrule
    Grafana     & \texttt{/etc/grafana/grafana.ini}              & Config generale, section \texttt{[smtp]} \\
    Prometheus  & \texttt{/etc/prometheus/prometheus.yml}         & Cibles de scraping \\
    Loki        & \texttt{/etc/loki/local-config.yaml}            & Stockage et retention des logs \\
    InfluxDB    & \texttt{/etc/influxdb/config.toml}              & Ports, auth, organisation \\
    \bottomrule
  \end{tabular}
\end{table}

\subsection{Alertmanager et notifications}

\begin{warnbox}
  \textbf{Alertmanager~: Configure mais canaux de notification a corriger.}\\
  Alertmanager est deploye et configure. Cependant, la configuration SMTP actuelle
  (section \texttt{[smtp]} dans \texttt{/etc/grafana/grafana.ini}) utilise un ancien
  compte Google avec un \textit{App Password} qui n'est plus valide. Pour que les
  alertes par e-mail fonctionnent, il faut soit~:\\[0.3em]
  \begin{itemize}[noitemsep]
    \item Remplacer le compte Google par un serveur SMTP dedie ou un relais SMTP
          valide (ex.~: Postfix local, SendGrid, Mailgun).
    \item Ou configurer un canal de notification alternatif directement dans Grafana
          (Grafana supporte Slack, Telegram, webhook, PagerDuty et d'autres, chacun
          necessitant l'ajout d'un \textit{contact point} dedie dans
          \texttt{Alerting > Contact points}).
  \end{itemize}
  Les regles d'alerte (seuils CPU, perte d'OSD Ceph, perte de quorum) sont a
  definir dans \texttt{Alerting > Alert rules} une fois le canal de notification valide.
\end{warnbox}

%============================================================
\section{GitLab et CI/CD Ansible}
%============================================================

\subsection{Instance GitLab}

Une instance GitLab est deployee sur l'infrastructure sous forme de
\textbf{conteneur LXC Turnkey} (image preconstruite \textit{Turnkey GitLab} pour Proxmox).
Elle heberge le depot du projet Ansible (\texttt{ansible-infra}).
Identifiants~: \texttt{root:ZakiIlyes7}.

\begin{successbox}
  \textbf{Statut~: Fonctionnel dans les conditions normales.}\\
  GitLab operationnel, depot \texttt{ansible-infra} disponible, runners configures.
\end{successbox}

\subsection{Probleme connu: Arret apres migration ou basculement HA}

\begin{warnbox}
  \textbf{GitLab peut s'arreter apres une migration a chaud ou un basculement HA.}\\
  Apres chaque deplacement du conteneur GitLab (declenche par CLBS ou par le
  mecanisme HA), les services GitLab peuvent ne pas redemarrer automatiquement.
  Pour relancer GitLab manuellement, entrer dans le conteneur puis executer~:
\end{warnbox}

\begin{lstlisting}[style=shellstyle]
gitlab-ctl start        # demarre tous les services GitLab
gitlab-ctl status       # verifie l'etat de chaque service
\end{lstlisting}

Si certains services restent en echec~:

\begin{lstlisting}[style=shellstyle]
gitlab-ctl restart <nom-du-service>   # ex. : gitlab-ctl restart puma
\end{lstlisting}

\subsection{Versions et mises a jour necessaires}

Le conteneur Turnkey GitLab presente deux problemes de version qu'il faudra
corriger manuellement (Turnkey ne fournit pas de guide de mise a jour pour ces deux points)~:

\begin{table}[H]
  \caption{Etat des versions du conteneur GitLab Turnkey}
  \centering
  \begin{tabular}{L{3.5cm}C{2.5cm}C{2.5cm}L{4cm}}
    \toprule
    \textbf{Composant} & \textbf{Version actuelle} & \textbf{Version cible} & \textbf{Methode} \\
    \midrule
    Debian (OS)  & 12 (Bookworm) & 13 (Trixie) & Mise a jour manuelle via \texttt{apt} \\
    GitLab CE    & 17.x          & 18.x        & Mise a jour manuelle via paquets GitLab \\
    \bottomrule
  \end{tabular}
\end{table}

\begin{warnbox}
  \textbf{Attention avant toute mise a jour~:} Sauvegarder l'etat du conteneur GitLab
  via Proxmox Backup Server avant d'entreprendre une mise a jour majeure.
  La mise a jour de Debian 12 vers 13 peut affecter les dependances de GitLab.
  Il est conseille de d'abord mettre a jour GitLab vers la derniere version 17.x stable,
  puis de passer a Debian 13, puis enfin de migrer vers GitLab 18.x.
\end{warnbox}

\subsection{Runner GitLab CI/CD: Probleme de linter}

Un \textit{runner} GitLab CI/CD est configure pour executer automatiquement
les verifications (\textit{lint}) et les controles syntaxiques (\textit{syntax-check})
sur les \textit{playbooks} Ansible a chaque \textit{push}.

\begin{warnbox}
  \textbf{Pipeline CI/CD~: Probleme au stage \texttt{lint}.}\\
  L'etape \texttt{ansible-lint} du pipeline echoue en raison de regles de linter
  trop strictes ou incompatibles avec les \textit{playbooks} actuels. Deux options
  pour resoudre ce probleme~:
  \begin{enumerate}[noitemsep]
    \item \textbf{Assouplir les regles}~: modifier \texttt{.ansible-lint} pour exclure
          les regles problematiques (ajouter \texttt{skip\_list:} avec les identifiants
          des regles qui echouent).
    \item \textbf{Supprimer l'etape lint}~: retirer les jobs \texttt{yaml-lint} et
          \texttt{ansible-lint} du fichier \texttt{.gitlab-ci.yml} et ne conserver que
          l'etape \texttt{syntax-check}, qui fonctionne correctement.
  \end{enumerate}
\end{warnbox}

Le fichier CI se trouve a \texttt{codes/ansible-infra/.gitlab-ci.yml}.
La configuration du linter est dans \texttt{.ansible-lint} et \texttt{.yamllint.yml}.

%============================================================
\section{Automatisation Ansible: Etat des playbooks}
%============================================================

\subsection{Playbooks disponibles}

Les \textit{playbooks} Ansible sont regroupes dans \texttt{codes/ansible-infra/playbooks/}.
L'interface graphique \textbf{Semaphore UI} (identifiants~: \texttt{admin:admin})
permet de les lancer sans passer par la ligne de commande.

\begin{table}[H]
  \caption{Etat des playbooks Ansible}
  \centering
  \begin{tabular}{L{4.5cm}L{3cm}L{6cm}}
    \toprule
    \textbf{Playbook} & \textbf{Statut} & \textbf{Remarques} \\
    \midrule
    \texttt{bootstrap.yml}   & \textcolor{vert}{\textbf{Fonctionnel}}  & Enrolement d'une nouvelle VM dans l'inventaire Ansible \\
    \texttt{hardening.yml}   & \textcolor{vert}{\textbf{Fonctionnel}}  & Durcissement SSH, mise a jour, configuration de base \\
    \texttt{ntp\_sync.yml}   & \textcolor{vert}{\textbf{Fonctionnel}}  & Synchronisation NTP des noeuds \\
    \texttt{pbs.yml}         & \textcolor{vert}{\textbf{Fonctionnel}}  & Configuration du Proxmox Backup Server \\
    \texttt{provision\_vm.yml} & \textcolor{orange}{\textbf{Partiel}} & Provisionnement unitaire de VM~; certains cas d'echec non resolus \\
    \texttt{provision\_vm\_batch.yml} & \textcolor{orange}{\textbf{Partiel}} & Provisionnement en lot~; pas entierement valide \\
    \bottomrule
  \end{tabular}
\end{table}

\begin{warnbox}
  \textbf{Playbooks de provisionnement~: Tests incomplets.}\\
  Les \textit{playbooks} \texttt{provision\_vm.yml} et \texttt{provision\_vm\_batch.yml}
  ont ete testes partiellement et ont echoue dans certains cas.
  L'integration d'un nouveau noeud au cluster via Ansible et la migration de VMs
  entre environnements (dev vers prod) n'ont pas ete entierement validees.
  Ces \textit{playbooks} necessitent des corrections avant utilisation en production.
\end{warnbox}

%============================================================
\section{Scalabilite: PXE / TFTP}
%============================================================

\subsection{Contexte}

Le deploiement de nouveaux noeuds physiques par demarrage reseau (PXE) a ete
envisage comme mecanisme de scalabilite horizontale de l'infrastructure.

\begin{warnbox}
  \textbf{PXE/TFTP~: Non deploye sur l'infrastructure reelle.}\\
  Aucun serveur PXE/TFTP n'a ete installe sur les serveurs physiques de l'infrastructure.
  Les tests realises l'ont ete \textbf{localement} sur les machines des developpeurs,
  en utilisant \textbf{iPXE} et un serveur TFTP. Ces tests ont rencontre des problemes
  lors du demarrage iPXE (difficultes de chainage du \textit{bootloader} et de
  chargement des images). La fonctionnalite reste en phase de preuve de concept
  et n'est pas prete pour un deploiement reel.
\end{warnbox}

Cette fonctionnalite est documentee dans le memoire (chapitre III, section Conception,
et perspectives) mais est explicitement marquee comme non realisee.

%============================================================
\section{Recapitulatif de l'etat de l'infrastructure}
%============================================================

\begin{table}[H]
  \caption{Etat global des composants}
  \centering
  \begin{tabular}{L{5.5cm}C{3cm}L{4.5cm}}
    \toprule
    \textbf{Composant} & \textbf{Statut} & \textbf{Remarques} \\
    \midrule
    Cluster Proxmox VE (3 noeuds) & \textcolor{vert}{\textbf{OK}} & Pleinement operationnel \\
    Reseau OVS + VxLAN            & \textcolor{vert}{\textbf{OK}} & 9 reseaux actifs \\
    Stockage Ceph                 & \textcolor{vert}{\textbf{OK}} & \texttt{HEALTH\_OK} \\
    Haute disponibilite (HA)      & \textcolor{vert}{\textbf{OK}} & Valide par test \\
    CLBS (DRS)                    & \textcolor{vert}{\textbf{OK}} & Migrations automatiques actives \\
    Application Sona-Web          & \textcolor{vert}{\textbf{OK}} & Flask + HAProxy + PostgreSQL \\
    Supervision (Grafana / InfluxDB) & \textcolor{vert}{\textbf{OK}} & Metriques disponibles \\
    GitLab (fonctionnement)       & \textcolor{vert}{\textbf{OK}} & Instance operationnelle \\
    Pare-feu Proxmox              & \textcolor{vert}{\textbf{OK}} & Active, mode DROP \\
    NAT OPNsense                  & \textcolor{vert}{\textbf{OK}} & Internet accessible depuis VMs \\
    Routage inter-zones OPNsense  & \textcolor{vert}{\textbf{OK}} & PROD / TEST / DMZ \\
    DHCP OPNsense                 & \textcolor{jaune}{\textbf{Partiel}} & Conflit avec switch physique \\
    Pare-feu OPNsense (inter-zones) & \textcolor{jaune}{\textbf{Partiel}} & Cas limites non testes \\
    DNS OPNsense                  & \textcolor{jaune}{\textbf{Partiel}} & Active, aucune regle personnalisee \\
    Pipeline CI/CD GitLab Ansible & \textcolor{jaune}{\textbf{Partiel}} & Etape lint en echec \\
    Playbooks de provisionnement  & \textcolor{jaune}{\textbf{Partiel}} & Cas d'echec non resolus \\
    Alertmanager / Notifications  & \textcolor{jaune}{\textbf{Partiel}} & SMTP invalide, a reconfigurer \\
    GitLab (versions)             & \textcolor{jaune}{\textbf{A mettre a jour}} & Debian 12, GitLab 17 \\
    Logs centralises (Loki)       & \textcolor{jaune}{\textbf{Partiel}} & Service deploye, sans sources de logs \\
    PXE / TFTP                    & \textcolor{rouge}{\textbf{Non deploye}} & Test local uniquement, problemes iPXE \\
    \bottomrule
  \end{tabular}
\end{table}

%============================================================
\section{Points d'action prioritaires}
%============================================================

\begin{enumerate}
  \item \textbf{Pipeline CI/CD}~: Corriger ou supprimer l'etape \texttt{ansible-lint}
        dans \texttt{.gitlab-ci.yml} pour restaurer un pipeline fonctionnel.

  \item \textbf{DHCP}~: Desactiver le DHCP sur le commutateur physique ou definir
        des adresses statiques sur toutes les VMs pour eliminer le conflit.

  \item \textbf{Ports applicatifs}~: Documenter et ouvrir au fur et a mesure les
        ports necessaires aux services en developpement (ex.~: 8000, 5432, 5000, 5423\ldots)
        dans le pare-feu Proxmox.

  \item \textbf{Alertmanager / SMTP}~: Remplacer la configuration SMTP invalide
        dans \texttt{/etc/grafana/grafana.ini} par un relais SMTP fonctionnel
        ou configurer un canal alternatif dans Grafana (Telegram, Slack, webhook).

  \item \textbf{GitLab: Mises a jour}~: Effectuer les mises a jour manuelles
        de Debian 12 vers 13 et de GitLab 17 vers 18 (sauvegarder le conteneur avant).

  \item \textbf{Playbooks de provisionnement}~: Identifier et corriger les cas
        d'echec de \texttt{provision\_vm.yml} et \texttt{provision\_vm\_batch.yml}
        avant toute utilisation en conditions reelles.

  \item \textbf{DNS OPNsense}~: Definir des enregistrements DNS internes pour
        les VMs afin de ne pas dependre des adresses IP brutes.

  \item \textbf{Migration vers CRS}~: Evaluer la mise a jour de Proxmox et le
        passage du CLBS au CRS natif pour un ordonnancement plus robuste.

  \item \textbf{Loki}~: Configurer les sources de logs (systemd, conteneurs, VMs)
        pour alimenter Loki et rendre la centralisation des logs effective.
\end{enumerate}

\end{document}
